\subsection{在 ad h 下稳定的子代数}

引理 1. 设 \(\mathrm{V}\) 是 \(\mathfrak{g}\) 的向量子空间，\(\mathrm{R(V)}\) 是 \(\alpha \in \mathrm{R}\) 中使 \(\mathfrak{g}^{\alpha} \subset \mathrm{V}\) 者组成的集合。则 \((\mathrm{V} \cap \mathfrak{h}) + \sum_{\alpha \in \mathrm{R(V)}} \mathfrak{g}^{\alpha}\) 是 \(\mathrm{V}\) 的在 \(\operatorname{ad} \mathfrak{h}\) 下稳定的最大向量子空间。

V 的向量子空间 W 在 ad h 下稳定当且仅当

\[
W = (W \cap \mathfrak {h}) + \sum_ {\alpha \in R} (W \cap \mathfrak {g} ^ {\alpha})
\]

（代数 第七章 §2 第 2 节 定理 1 之推论 1）。于是 V 的在 ad \(\mathfrak{h}\) 下稳定的最大向量子空间是 \((\mathrm{V} \cap \mathfrak{h}) + \sum_{\alpha \in \mathrm{R}} (\mathrm{V} \cap \mathfrak{g}^{\alpha})\) 。但 \(\mathrm{V} \cap \mathfrak{g}^{\alpha} = \mathfrak{g}^{\alpha}\) （对 \(\alpha \in \mathrm{R}(\mathrm{V})\) ），而 \(\mathrm{V} \cap \mathfrak{g}^{\alpha} = 0\) （对 \(\alpha \notin \mathrm{R}(\mathrm{V})\) ），因为 \(\dim \mathfrak{g}^{\alpha} = 1\) 。证毕。

对 R 的任意子集 P，置

\[
\mathfrak {g} ^ {\mathrm{P}} = \sum_ {\alpha \in \mathrm{P}} \mathfrak {g} ^ {\alpha} \quad \mathfrak {h} _ {\mathrm{P}} = \sum_ {\alpha \in \mathrm{P}} \mathfrak {h} _ {\alpha}.
\]

若 \(\mathrm{P}\subset \mathrm{R}\) 且 \(\mathbf{Q}\subset \mathbb{R}\) ，我们显然有

\[
\begin{array}{c} {[ \mathfrak {h}, \mathfrak {g} ^ {\mathrm{P}} ] \subset \mathfrak {g} ^ {\mathrm{P}}} \\ {[ \mathfrak {g} ^ {\mathrm{P}}, \mathfrak {g} ^ {\mathrm{Q}} ] = \mathfrak {g} ^ {(\mathrm{P+Q}) \cap \mathrm{R}} + \mathfrak {h} _ {\mathrm{P} \cap (- \mathrm{Q})}.} \end{array}\tag{1}
\]

\begin{enumerate}
\def\labelenumi{(\arabic{enumi})}
\setcounter{enumi}{1}
\tightlist
\item
\end{enumerate}

回顾（第六章 §1 第 7 节 定义 4），R 的子集 P 称为闭的，如果条件 \(\alpha \in P\) ，\(\beta \in P\) ，\(\alpha + \beta \in R\) 蕴含 \(\alpha + \beta \in P\) ，换言之，如果 \((P + P) \cap R \subset P\) 。

引理 2. 设 \(\mathfrak{h}'\) 是 \(\mathfrak{h}\) 的向量子空间，\(\mathrm{P}\) 是 \(\mathrm{R}\) 的子集。则 \(\mathfrak{h}' + \mathfrak{g}^{\mathrm{P}}\) 是 \(\mathfrak{g}\) 的子代数当且仅当 \(\mathrm{P}\) 是 \(\mathrm{R}\) 的闭子集并且

\[
\mathfrak {h} ^ {\prime} \supset \mathfrak {h} _ {\mathrm{P} \cap (- \mathrm{P})}.
\]

事实上，

\[
[ \mathfrak {h} ^ {\prime} + \mathfrak {g} ^ {\mathrm{P}}, \mathfrak {h} ^ {\prime} + \mathfrak {g} ^ {\mathrm{P}} ] = [ \mathfrak {h} ^ {\prime}, \mathfrak {g} ^ {\mathrm{P}} ] + [ \mathfrak {g} ^ {\mathrm{P}}, \mathfrak {g} ^ {\mathrm{P}} ] = \mathfrak {h} _ {\mathrm{P} \cap (- \mathrm{P})} + [ \mathfrak {h} ^ {\prime}, \mathfrak {g} ^ {\mathrm{P}} ] + \mathfrak {g} ^ {(\mathrm{P} + \mathrm{P}) \cap \mathrm{R}}.
\]

因此 \(\mathfrak{h}' + \mathfrak{g}^{\mathrm{P}}\) 是 \(\mathfrak{g}\) 的子代数当且仅当

\[
\mathfrak {h} _ {\mathrm{P} \cap (- \mathrm{P})} \subset \mathfrak {h} ^ {\prime} \text { and } \mathfrak {g} ^ {(\mathrm{P} + \mathrm{P}) \cap \mathrm{R}} \subset \mathfrak {g} ^ {\mathrm{P}}
\]

这就证明了引理。

命题 1. (i) g 的在 ad h 下稳定的子代数正是形如 \(h' + g^{P}\) 的向量子空间，其中 P 是 R 的闭子集，而 \(h'\) 是 h 的包含 \(\mathfrak{h}_{\mathrm{P}\cap(-\mathrm{P})}\) 的向量子空间。

\begin{enumerate}
\def\labelenumi{(\roman{enumi})}
\setcounter{enumi}{1}
\tightlist
\item
  设 \(\mathfrak{h}'\) ，\(\mathfrak{h}''\) 是 \(\mathfrak{h}\) 的向量子空间，\(\mathrm{P},\mathrm{Q}\) 是 \(\mathrm{R}\) 的闭子集，且 \(\mathfrak{h}'\supset \mathfrak{h}_{\mathrm{P}\cap (-\mathrm{P})}\) ，\(\mathfrak{h}''\subset \mathfrak{h}'\) ，\(\mathrm{Q}\subset \mathrm{P}\) 。则 \(\mathfrak{h}'' + \mathfrak{g}^{\mathrm{Q}}\) 是 \(\mathfrak{h}' + \mathfrak{g}^{\mathrm{P}}\) 的理想当且仅当
\end{enumerate}

\[
(P + Q) \cap R \subset Q \text { and } \mathfrak {h} _ {P \cap (- Q)} \subset \mathfrak {h} ^ {\prime \prime} \subset \bigcap_ {\alpha \in P, \alpha \notin Q} \operatorname{Ker} \alpha .
\]

断言 (i) 由引理 1 与引理 2 立即得出。设 \(\mathfrak{h}',\mathfrak{h}''\) ，P, Q 如 (ii)。则

\[
[ \mathfrak {h} ^ {\prime} + \mathfrak {g} ^ {\mathrm{P}}, \mathfrak {h} ^ {\prime \prime} + \mathfrak {g} ^ {\mathrm{Q}} ] = \mathfrak {h} _ {\mathrm{P} \cap (- \mathrm{Q})} + [ \mathfrak {h} ^ {\prime}, \mathfrak {g} ^ {\mathrm{Q}} ] + [ \mathfrak {h} ^ {\prime \prime}, \mathfrak {g} ^ {\mathrm{P}} ] + \mathfrak {g} ^ {(\mathrm{P} + \mathrm{Q}) \cap \mathrm{R}}.
\]

于是，\(\mathfrak{h}'' + \mathfrak{g}^{\mathrm{Q}}\) 是 \(\mathfrak{h}' + \mathfrak{g}^{\mathrm{P}}\) 的理想当且仅当

\[
\mathfrak {h} _ {\mathrm{P} \cap (- \mathrm{Q})} \subset \mathfrak {h} ^ {\prime \prime}, [ \mathfrak {h} ^ {\prime \prime}, \mathfrak {g} ^ {\mathrm{P}} ] \subset \mathfrak {g} ^ {\mathrm{Q}}, \mathfrak {g} ^ {(\mathrm{P} + \mathrm{Q}) \cap \mathrm{R}} \subset \mathfrak {g} ^ {\mathrm{Q}}.
\]

这就蕴含 (ii)。

命题 2. 设 \(\mathfrak{a}\) 是 \(\mathfrak{g}\) 的在 ad \(\mathfrak{h}\) 下稳定的子代数，并设 \(\mathfrak{h}' \subset \mathfrak{h}\) ，\(\mathrm{P} \subset \mathrm{R}\) 使 \(\mathfrak{a} = \mathfrak{h}' + \mathfrak{g}^{\mathrm{P}}\) 。

\begin{enumerate}
\def\labelenumi{(\roman{enumi})}
\item
  设 \(\mathfrak{k}\) 是 \(x\in \mathfrak{h}'\) 中使 \(\alpha (x) = 0\) 对一切 \(\alpha \in \mathrm{P}\cap (-\mathrm{P})\) 成立者组成的集合。\(\mathfrak{a}\) 的根基是 \(\mathfrak{k} + \mathfrak{g}^{\mathrm{Q}}\) ，其中 \(\mathrm{Q}\) 是 \(\alpha \in \mathrm{P}\) 中使 \(-\alpha \notin \mathrm{P}\) 者组成的集合。此外，\(\mathfrak{g}^{\mathrm{Q}}\) 是 \(\mathfrak{a}\) 的一个幂零理想。
\item
  a 半单当且仅当 P = -P 且 \(h' = h_{P}\) 。
\item
  \(\mathfrak{a}\) 可解当且仅当 \(\mathrm{P} \cap (-\mathrm{P}) = \emptyset\) 。这时 \([\mathfrak{a}, \mathfrak{a}] = \mathfrak{g}^{\mathrm{S}}\) ，其中
\end{enumerate}

\[
\mathrm{S} = ((\mathrm{P} + \mathrm{P}) \cap \mathrm{R}) \cup \{\alpha \in \mathrm{P} | \alpha (\mathfrak {h} ^ {\prime}) \neq 0 \}.
\]

\begin{enumerate}
\def\labelenumi{(\roman{enumi})}
\setcounter{enumi}{3}
\item
  \(\mathfrak{a}\) 在 \(\mathfrak{g}\) 中约化当且仅当 \(\mathrm{P} = -\mathrm{P}\) 。
\item
  \(\mathfrak{a}\) 由幂零元素组成当且仅当 \(\mathfrak{h}' = 0\) 。这时 \(\mathrm{P} \cap (-\mathrm{P}) = \emptyset\) ，且 \(\mathfrak{a}\) 是幂零的。
\end{enumerate}

我们证明 (v)。若 \(\mathfrak{a}\) 由幂零元素组成，则 \(\mathfrak{a}\) 显然是幂零的，且 \(\mathfrak{h}' = 0\) ，因为 \(\mathfrak{h}\) 的元素都是半单的。设 \(\mathfrak{h}' = 0\) 。由命题 1 (i)，\(P \cap (-P) = \varnothing\) 。由第六章 §1 第 7 节 命题 22，存在 R 的一个房 C 使 \(P \subset R_{+}(C)\) 。因此，存在具有下列性质的整数 \(n > 0\) ：若 \(\alpha_{1}, \ldots, \alpha_{n} \in P\) 且 \(\beta \in R \cup \{0\}\) ，则

\[
\alpha_ {1} + \dots + \alpha_ {n} + \beta \notin \mathrm{R} \cup \{0 \}.
\]

这蕴含 \(\mathfrak{g}^{\mathrm{P}}\) 的每个元素都是幂零的，于是得到 (v)。

我们证明 (iii)。若 \(\mathrm{P} \cap (-\mathrm{P}) = \varnothing\) ，则 \(\mathfrak{g}^{\mathrm{P}}\) 是 \(\mathfrak{g}\) 的子代数（命题 1 (i)），且由 (v) 是幂零的。现在

\[
[ \mathfrak {a}, \mathfrak {a} ] = [ \mathfrak {h} ^ {\prime}, \mathfrak {g} ^ {\mathrm{P}} ] + [ \mathfrak {g} ^ {\mathrm{P}}, \mathfrak {g} ^ {\mathrm{P}} ] = [ \mathfrak {h} ^ {\prime}, \mathfrak {g} ^ {\mathrm{P}} ] + \mathfrak {g} ^ {(\mathrm{P} + \mathrm{P}) \cap \mathrm{R}} \subset \mathfrak {g} ^ {\mathrm{P}},
\]

于是 \(\mathfrak{a}\) 可解，且 \([\mathfrak{a},\mathfrak{a}]\) 由命题中的公式给出。若 \(\mathrm{P}\cap (-\mathrm{P})\neq \emptyset\) ，取 \(\alpha \in \mathrm{P}\) 使 \(-\alpha \in \mathrm{P}\) 。则 \(\mathfrak{h}_{\alpha} + \mathfrak{g}^{\alpha} + \mathfrak{g}^{-\alpha}\) 是 \(\mathfrak{a}\) 的一个单子代数，故 \(\mathfrak{a}\) 不可解。

我们证明 (i)。由于 P 是闭的，\((\mathrm{P} + \mathrm{Q}) \cap \mathrm{R} \subset \mathrm{P}\) 。若 \(\alpha \in P, \beta \in Q\) 且 \(\alpha + \beta \in R\) ，则不可能有 \(\alpha + \beta \in -P\) ，因为，P 既为闭的，这将蕴含 \(-\beta = -(\alpha + \beta) + \alpha \in P\) ，而 \(\beta \in Q\) ；于是，\((\mathrm{P} + \mathrm{Q}) \cap \mathrm{R} \subset \mathrm{Q}\) 。这证明了 \(g^{Q}\) 是 a 的理想，且由 (v) 是幂零的。我们有 \(\mathrm{P} \cap (-\mathrm{Q}) = \varnothing\) ，且 \(\mathrm{P} \cap (-\mathrm{P}) = \mathrm{P} \cap \complement Q\) ，故 \(\mathfrak{h}_{\mathrm{P} \cap (-\mathrm{Q})} \subset \mathfrak{k} \subset \bigcap_{\alpha \in \mathrm{P}, \alpha \notin Q} \operatorname{Ker} \alpha\) 。由命题 1 (ii)，

\(\mathfrak{k} + \mathfrak{g}^{\mathrm{Q}}\) 是 \(\mathfrak{a}\) 的理想。由于 \(\mathrm{Q} \cap (-\mathrm{Q}) = \emptyset\) ，由 (iii) 这个理想是可解的。因此它含于根基 \(\mathfrak{r}\) （\(\mathfrak{a}\) 的）。由于 \(\mathfrak{r}\) 在 \(\mathfrak{a}\) 的每个导子下稳定，\(\mathfrak{r}\) 在 ad \(\mathfrak{h}\) 下稳定。因此存在子集 \(S\) （\(P\) 的），使 \(\mathfrak{r} = (\mathfrak{r} \cap \mathfrak{h}) + \mathfrak{g}^{\mathrm{S}}\) 。设 \(\alpha \in S\) 且 \(-\alpha \in P\) 。则 \(\mathfrak{h}_{\alpha} = [\mathfrak{g}^{\alpha}, \mathfrak{g}^{-\alpha}] \subset \mathfrak{r}\) ，故 \(\mathfrak{g}^{-\alpha} = [\mathfrak{h}_{\alpha}, \mathfrak{g}^{-\alpha}] \subset \mathfrak{r} = 0\) ，从而 \(-\alpha \in S\) ；由 (iii)，这与 \(\mathfrak{r}\) 可解矛盾。因而 \(S \subset Q\) 。最后，若 \(x \in \mathfrak{r} \cap \mathfrak{h}\) 且 \(\alpha \in P \cap (-P)\) ，则 \([x, \mathfrak{g}^{\alpha}] \subset \mathfrak{g}^{\alpha} \cap \mathfrak{r} = 0\) ，故 \(\alpha(x) = 0\) ；这表明 \(x \in \mathfrak{k}\) 。于是 \(\mathfrak{r} \subset \mathfrak{k} + \mathfrak{g}^{\mathrm{Q}}\) ，(i) 的证明完毕。

我们证明 (iv)。由 (i)，\(\mathfrak{a}\) 在 \(\mathfrak{g}\) 上的伴随表示是半单的，当且仅当 \(\mathrm{ad}_{\mathfrak{g}}x\) 对一切 \(x\in \mathfrak{k} + \mathfrak{g}^{\mathrm{Q}}\) 是半单的（第一章 §6 第 5 节 定理 4）；由 (v)，此事成立当且仅当 \(\mathrm{Q} = \varnothing\) ，换言之 \(\mathrm{P} = -\mathrm{P}\) 。

我们证明 (ii)。若 \(\mathfrak{a}\) 半单，则由 (i) 有 \(P = -P\) ，故 \(\mathfrak{h}_{\mathrm{P}} \subset \mathfrak{h}'\) ；再者，\(\mathfrak{a} = [\mathfrak{a}, \mathfrak{a}] \subset \mathfrak{h}_{\mathrm{P}} + \mathfrak{g}^{\mathrm{P}}\) ，从而 \(\mathfrak{h}' = \mathfrak{h}_{\mathrm{P}}\) 。若 \(P = -P\) 且 \(\mathfrak{h}' = \mathfrak{h}_{\mathrm{P}}\) ，则 \(\mathfrak{a}\) 由 (iv) 是约化的，且 \(\mathfrak{a} = \sum_{\alpha \in P} \mathfrak{s}_{\alpha}\) ，故 \(\mathfrak{a} = [\mathfrak{a}, \mathfrak{a}]\) ，\(\mathfrak{a}\) 是半单的。

命题 3. 设 \(\mathfrak{a}\) 是 \(\mathfrak{g}\) 的在 \(\operatorname{ad}(\mathfrak{h})\) 下稳定的半单子代数，并设 \(\mathrm{P}\) 是 \(\mathrm{R}\) 的子集，它使 \(\mathfrak{a} = \mathfrak{h}_{\mathrm{P}} + \mathfrak{g}^{\mathrm{P}}\) 。

\begin{enumerate}
\def\labelenumi{(\roman{enumi})}
\item
  \(\mathfrak{h}_{\mathrm{P}}\) 是 \(\mathfrak{a}\) 的一个分裂嘉当子代数。
\item
  \((\mathfrak{a},\mathfrak{h}_{\mathrm{P}})\) 的根系是诸元素到 \(\mathfrak{h}_{\mathrm{P}}\) 上的限制所成的集合，这些元素取自 \(\mathrm{P}\) 。
\end{enumerate}

由于 \(h_{P}\) 在 ad h 下稳定，它在 a 中的正规化子在 ad h 下稳定，因而具有形式 \(h_{P} + g^{Q}\) ，其中 \(Q \subset P\) （引理 1）。若 \(\alpha \in Q\) ，

\[
\mathfrak {g} ^ {\alpha} = \left[ \mathfrak {h} _ {\alpha}, \mathfrak {g} ^ {\alpha} \right] \subset \left[ \mathfrak {h} _ {\mathrm{P}}, \mathfrak {g} ^ {\alpha} \right] \subset \mathfrak {h} _ {\mathrm{P}},
\]

这是荒谬的。于是 \(Q = \varnothing\) ，\(h_{P}\) 是它在 a 中自己的正规化子。这证明了 \(h_{P}\) 是 a 的一个嘉当子代数。若 \(x \in h_{P}\) ，则 \(ad_{g}x\) ，从而更有 \(ad_{a}x\) ，是可三角化的。于是 (i) 得证，而 (ii) 是明显的。

由命题 1 (i)，g 的包含 h 的子代数正是那些集合 \(h + g^{P}\) ，其中 P 是 R 的闭子集。由第七章 §3 命题 3，\(h + g^{P}\) 的每个嘉当子代数都是 g 的嘉当子代数。

命题 4. 设 \(\mathfrak{a}\) 是 \(\mathfrak{g}\) 的包含 \(\mathfrak{h}\) 的子代数，\(x\) 是 \(\mathfrak{a}\) 的元素，\(s\) 与 \(n\) 是它的半单分量与幂零分量。则 \(s \in \mathfrak{a}\) 且 \(n \in \mathfrak{a}\) 。

我们有 \((\operatorname{ad} x)\mathfrak{a} \subset \mathfrak{a}\) ，故 \((\operatorname{ad} s)\mathfrak{a} \subset \mathfrak{a}\) 且 \((\operatorname{ad} n)\mathfrak{a} \subset \mathfrak{a}\) 。由于 a 是它在 g 中自己的正规化子（第七章 §2 第 1 节 命题 4 之推论 4），\(s \in a\) 且 \(n \in a\) 。

命题 5. 设 P 是 R 的闭子集。

\begin{enumerate}
\def\labelenumi{(\roman{enumi})}
\item
  \(\mathfrak{h} + \mathfrak{g}^{\mathrm{P}}\) 可解当且仅当 \(\mathrm{P} \cap (-\mathrm{P}) = \varnothing\) 。这时，\([\mathfrak{h} + \mathfrak{g}^{\mathrm{P}}, \mathfrak{h} + \mathfrak{g}^{\mathrm{P}}] = \mathfrak{g}^{\mathrm{P}}\) 。
\item
  \(\mathfrak{h} + \mathfrak{g}^{\mathrm{P}}\) 约化当且仅当 \(\mathrm{P} = -\mathrm{P}\) 。
\end{enumerate}

断言 (i) 由命题 2 (iii) 得出。若 \(\mathrm{P} = -\mathrm{P}\) ，则 \(\mathfrak{h} + \mathfrak{g}^{\mathrm{P}}\) 是约化的（命题 2 (iv)）。设 \(\mathfrak{a} = \mathfrak{h} + \mathfrak{g}^{\mathrm{P}}\) 是约化的。则

\[
\mathfrak {g} ^ {\mathrm{P}} = \left[ \mathfrak {h}, \mathfrak {g} ^ {\mathrm{P}} \right] \subset \left[ \mathfrak {a}, \mathfrak {a} \right] \subset \mathfrak {h} + \mathfrak {g} ^ {\mathrm{P}},
\]

于是 \([a, a]\) 具有形式 \(h' + g^{P}\) ，其中 \(h' \subset h\) ；由于 \([a, a]\) 半单，故 P = -P （命题 2 (ii)）。

